\documentclass[11pt,a4paper]{article}

\usepackage[margin=1in]{geometry}
\usepackage{amsmath}
\usepackage{amssymb}
\usepackage{booktabs}
\usepackage[hidelinks]{hyperref}

\newcommand{\Reals}{\mathbb{R}}
\newcommand{\SO}{\mathrm{SO}}
\newcommand{\SE}{\mathrm{SE}}
\newcommand{\norm}[1]{\left\lVert #1 \right\rVert}
\newcommand{\Tf}[2]{{}^{#1}\mathbf{T}_{#2}}
\newcommand{\Rot}[2]{{}^{#1}\mathbf{R}_{#2}}

\title{Navigation and Perception: the mathematics}
\author{HSR search-and-fetch pipeline}
\date{}

\begin{document}
\maketitle

\begin{abstract}
This note states, in full, the mathematics the navigation and perception halves
of the pipeline actually implement: how a camera image becomes a point in the
map, and how a point in the map becomes a place for the robot to stand. Every
formula corresponds to code in \texttt{core/}, and the file and function are
named beside each one. Terms are defined as they first appear, so the note can
be read without prior robotics background.
\end{abstract}

\tableofcontents

%======================================================================
\section{Notation and conventions}

\subsection{Points, frames, and why frames matter}

A \emph{point} is a location in space, written $\mathbf{p} = (x,y,z)^\top \in \Reals^3$.
Three numbers alone are meaningless without saying \emph{what they are measured
from}. That reference is a \emph{frame}: an origin plus three mutually
perpendicular unit axes. The same physical object has different coordinates in
different frames.

Four frames appear in this pipeline.

\begin{table}[h]
\centering
\begin{tabular}{lll}
\toprule
Frame & Origin & Used by \\
\midrule
\texttt{map}       & Fixed corner of the environment       & Nav2 goals, the scene graph \\
\texttt{odom}      & Where the robot booted                & MoveIt planning, saved grasps \\
\texttt{base\_footprint} & The robot, on the floor         & The robot's own pose \\
camera optical     & The camera lens                       & Raw depth measurements \\
\bottomrule
\end{tabular}
\end{table}

\texttt{map} and \texttt{odom} coincide at start-up and drift apart as the
localiser (AMCL) corrects the robot's estimated position. Mixing them is a real
source of error: \S\ref{sec:frames} returns to this.

We write ${}^{A}\mathbf{p}$ for a point expressed in frame $A$.

\subsection{Rotations}

A \emph{rotation matrix} $\mathbf{R} \in \SO(3)$ is a $3\times3$ matrix whose
columns are unit vectors at right angles to each other:
\begin{equation}
\mathbf{R}^\top \mathbf{R} = \mathbf{I}, \qquad \det \mathbf{R} = +1 .
\end{equation}
The first condition says the axes stay unit-length and perpendicular; the second
rules out mirror images. Two consequences are used constantly:
$\mathbf{R}^{-1} = \mathbf{R}^\top$ (inverting a rotation is free), and
$\norm{\mathbf{R}\mathbf{v}} = \norm{\mathbf{v}}$ (rotations preserve lengths).

Column $j$ of $\mathbf{R}$ is the image of the $j$-th unit axis. Writing
$\mathbf{e}_1=(1,0,0)^\top$, $\mathbf{e}_2=(0,1,0)^\top$,
$\mathbf{e}_3=(0,0,1)^\top$, the third column $\mathbf{R}\mathbf{e}_3$ is where
the frame's own $+z$ axis points. This single fact is what lets the code extract
a gripper's approach direction from a pose without any library call.

\subsection{Homogeneous transforms}

To move a point between frames we must rotate \emph{and} translate. Doing both at
once is convenient if we append a $1$ to every point, giving a
\emph{homogeneous} coordinate $(x,y,z,1)^\top$, and use a $4\times4$ matrix:
\begin{equation}
\Tf{A}{B} =
\begin{bmatrix}
\Rot{A}{B} & {}^{A}\mathbf{t}_{B} \\
\mathbf{0}^\top & 1
\end{bmatrix} \in \SE(3),
\qquad
\begin{bmatrix} {}^{A}\mathbf{p} \\ 1 \end{bmatrix}
= \Tf{A}{B}
\begin{bmatrix} {}^{B}\mathbf{p} \\ 1 \end{bmatrix}.
\label{eq:homog}
\end{equation}
Read $\Tf{A}{B}$ as ``the pose of frame $B$, expressed in frame $A$'', or
equivalently ``the map that takes $B$-coordinates to $A$-coordinates''. Here
$\Rot{A}{B}$ is the rotation and ${}^{A}\mathbf{t}_{B}$ the position of $B$'s
origin in $A$.

Transforms compose by matrix multiplication, and the adjacent subscripts cancel:
\begin{equation}
\Tf{A}{C} = \Tf{A}{B}\, \Tf{B}{C}.
\label{eq:compose}
\end{equation}
Inversion is closed-form, needing no general matrix inverse:
\begin{equation}
\Tf{A}{B}^{-1} = \Tf{B}{A} =
\begin{bmatrix}
\Rot{A}{B}^\top & -\Rot{A}{B}^\top\, {}^{A}\mathbf{t}_{B} \\
\mathbf{0}^\top & 1
\end{bmatrix}.
\end{equation}

A $4\times4$ matrix of this form is used for two different jobs, and it is worth
keeping them apart: it can be a \emph{change of coordinates} (the same point,
described from a different frame) or a \emph{pose} (a physical object's position
and orientation). The algebra is identical; the interpretation is not.

%======================================================================
\section{Perception}

The perception half answers: \emph{given an image and a word, where is that
object in the map?}

\subsection{The pinhole camera model}

A depth camera returns, for each pixel $(u,v)$, a distance $d(u,v)$ along the
camera's viewing axis. To turn a pixel into a 3D point we need the camera's
\emph{intrinsic matrix}
\begin{equation}
\mathbf{K} =
\begin{bmatrix}
f_x & 0 & c_x \\
0 & f_y & c_y \\
0 & 0 & 1
\end{bmatrix},
\end{equation}
whose four meaningful entries are:
\begin{itemize}
\item $f_x, f_y$ --- the \emph{focal lengths} in pixels. Larger values mean a
      narrower field of view: the same object spans more pixels.
\item $c_x, c_y$ --- the \emph{principal point}, the pixel where the optical axis
      meets the image. Usually near the image centre but rarely exactly there.
\end{itemize}

The pinhole model states that a point ${}^{C}\mathbf{p} = (X,Y,Z)^\top$ in the
camera frame projects to the pixel
\begin{equation}
u = f_x \frac{X}{Z} + c_x, \qquad v = f_y \frac{Y}{Z} + c_y .
\end{equation}
Objects twice as far away appear half as large: that is the division by $Z$.

\subsection{Back-projection}

Projection loses one dimension, so it cannot be undone from a pixel alone. The
depth channel restores exactly the missing number. Setting $Z = d(u,v)$ and
solving gives \emph{back-projection}:
\begin{equation}
\boxed{\;
{}^{C}\mathbf{p}(u,v) =
\left(
\frac{(u - c_x)\, d(u,v)}{f_x},\;
\frac{(v - c_y)\, d(u,v)}{f_y},\;
d(u,v)
\right)^\top .
\;}
\label{eq:deproject}
\end{equation}
\textit{Implemented in} \texttt{core/perception/pointcloud.py}, \texttt{deproject}.

Only pixels with a valid return are used. Writing $M \subseteq \{1..H\}\times\{1..W\}$
for the set of pixels the segmenter marked as belonging to the object, the
retained set is
\begin{equation}
V = \{\, (u,v) \in M \;:\; d(u,v) > 0 \ \text{and}\ d(u,v) \ \text{finite} \,\}.
\end{equation}
Zero depth is not a measurement at the origin; it is the sensor's code for ``no
return''. Treating it as a point would drag every later average towards the
camera.

The camera reports millimetres as 16-bit integers on the real robot and metres as
floats in simulation, so a unit conversion precedes all of this. All mathematics
below is in metres.

\subsection{Segmentation and the object point set}

An open-vocabulary segmenter (SAM3) maps an image and a text prompt to a set of
candidate instance masks with confidence scores,
$\{(M_i, s_i)\}_{i=1}^{n}$. The pipeline keeps the most confident one,
\begin{equation}
i^\star = \arg\max_i s_i, \qquad M = M_{i^\star},
\end{equation}
and treats $n = 0$ as ``the object is not visible from here'' --- a signal the
search loop acts on, not an error.
\textit{Implemented in} \texttt{core/perception/sam3\_client.py}.

Applying \eqref{eq:deproject} over $V$ yields the object's point set in camera
coordinates,
\begin{equation}
\mathcal{P}_C = \{\, {}^{C}\mathbf{p}(u,v) \;:\; (u,v) \in V \,\} \subset \Reals^3 .
\end{equation}

\subsection{Into a world frame}
\label{sec:frames}

The camera moves with the head, so camera coordinates are useless for planning.
Using \eqref{eq:homog} with the transform read from the robot's TF tree,
\begin{equation}
\mathcal{P}_W = \{\, \Tf{W}{C}\, \mathbf{p} \;:\; \mathbf{p} \in \mathcal{P}_C \,\},
\end{equation}
where $W$ is the target frame. In code this is done without homogeneous padding,
using the equivalent form
\begin{equation}
{}^{W}\mathbf{p} = \Rot{W}{C}\, {}^{C}\mathbf{p} + {}^{W}\mathbf{t}_{C},
\end{equation}
applied to all points at once as $\mathbf{P}\,\Rot{W}{C}^\top + \mathbf{t}^\top$
for a stacked $N \times 3$ array $\mathbf{P}$.
\textit{Implemented in} \texttt{core/perception/pointcloud.py}, \texttt{transform\_points}.

\paragraph{Which $W$.} The choice is not cosmetic. Anything handed to Nav2 must be
in \texttt{map}, because that is the navigation stack's global frame; anything
handed to MoveIt is in \texttt{odom}. The two agree at start-up and separate as
the localiser corrects, so a quantity computed in one and consumed in the other
degrades silently over a run rather than failing outright.

\subsection{The object's position}

A single position summarising the object is the arithmetic mean, or
\emph{centroid}, of its points:
\begin{equation}
\boxed{\;
\bar{\mathbf{p}} = \frac{1}{|\mathcal{P}_W|} \sum_{\mathbf{p} \in \mathcal{P}_W} \mathbf{p}.
\;}
\label{eq:centroid}
\end{equation}

Two properties are worth stating because they are relied on. First, the centroid
commutes with rigid transforms,
\begin{equation}
\overline{\Tf{W}{C}\,\mathcal{P}_C} = \Tf{W}{C}\, \overline{\mathcal{P}_C},
\end{equation}
so it does not matter whether one averages before or after changing frames.
Second, it is \emph{not} robust: a single stray point at $10$~m --- a depth
outlier at the mask's edge, where the mask overlaps the background --- shifts the
mean by $10/|\mathcal{P}_W|$ metres. With a few thousand points that is
millimetres, which is why the plain mean suffices here; with a ragged mask over
few points it would not, and a median or a trimmed mean would be the fix.

An axis-aligned bounding box is also recorded, as the componentwise extremes
\begin{equation}
\mathbf{b}^{-} = \min_{\mathbf{p} \in \mathcal{P}_W} \mathbf{p}, \qquad
\mathbf{b}^{+} = \max_{\mathbf{p} \in \mathcal{P}_W} \mathbf{p},
\end{equation}
so a consumer can register the object as a collision obstacle.

\subsection{Grasp poses and the gripper frame convention}

The grasp generator returns poses $\mathbf{G}_i \in \SE(3)$ in its own
\emph{canonical grasp frame}, defined by the convention
\begin{equation}
+z \ \text{is the approach direction}, \qquad
+x \ \text{is the closing direction}.
\end{equation}
The robot's hand link obeys a different convention: for the HSRC hand, forward
kinematics gives $+z$ as the approach direction but $+y$ as the closing
direction. The two differ by a rotation of $-90^\circ$ about the shared approach
axis:
\begin{equation}
\mathbf{R}_{\text{base}} =
\begin{bmatrix}
0 & 1 & 0 \\
-1 & 0 & 0 \\
0 & 0 & 1
\end{bmatrix}
= \mathbf{R}_z(-\tfrac{\pi}{2}).
\end{equation}
This matrix maps hand-link coordinates into canonical coordinates. To express a
grasp as a pose \emph{of the hand link}, compose on the right:
\begin{equation}
\boxed{\;
\mathbf{G}^{\text{hand}} = \mathbf{G}^{\text{canonical}} \; \mathbf{R}_{\text{base}} .
\;}
\label{eq:baserot}
\end{equation}
\textit{Implemented in} \texttt{core/grasping/graspgenx\_client.py}.

The ordering follows from \eqref{eq:compose}: with $H$ the hand-link frame and
$G$ the canonical grasp frame, $\Tf{W}{H} = \Tf{W}{G}\,\Tf{G}{H}$ and
$\Tf{G}{H} = \mathbf{R}_{\text{base}}$.

Because $\mathbf{R}_{\text{base}}$ is a rotation \emph{about} $z$, it leaves the
approach axis fixed. Omitting it therefore leaves every approach-direction
computation correct while rotating the jaws $90^\circ$ --- an error invisible in
every intermediate check and fatal only at contact, where the fingers close
across the wrong axis of the object.

Finally, since the generator is trained on object-centred clouds, the cloud is
centred before the query and the offset restored afterwards. Only the
translation part is affected:
\begin{equation}
\mathcal{P}' = \{\mathbf{p} - \bar{\mathbf{p}}\}, \qquad
\mathbf{G}_i \leftarrow \mathbf{G}_i + \bar{\mathbf{p}}\,\mathbf{e}_4^\top,
\end{equation}
which adds $\bar{\mathbf{p}}$ to the last column of each pose and leaves the
rotation block untouched. Rotations are unaffected because a translation of the
input translates the grasp without turning it.

%======================================================================
\section{Navigation}

The navigation half answers two distinct questions that are easy to conflate:
\emph{where can the robot stand to \textbf{see} a piece of furniture}, and
\emph{where can it stand so the \textbf{arm} reaches a given hand pose}. They
have different answers, and \S\ref{sec:disjoint} shows the two never coincide.

\subsection{Standing off a piece of furniture}

Furniture is modelled as an axis-aligned box of full dimensions
$(D_x, D_y, D_z)$ centred at $\mathbf{c}$. Write the half-extents
$a = D_x/2$, $b = D_y/2$.

We want a ring of poses clear of the furniture's \emph{footprint} --- its
shadow on the floor --- rather than a circle of fixed radius about the centroid.
For a non-square piece these differ substantially.

How far the footprint itself extends in direction
$\mathbf{u}(\theta) = (\cos\theta, \sin\theta)^\top$ is given by its
\emph{support function}: the greatest projection of any point of the shape onto
that direction,
\begin{equation}
h(\theta) = \max_{\mathbf{q} \in \text{box}} \mathbf{q}^\top \mathbf{u}(\theta).
\end{equation}
For an axis-aligned rectangle centred at the origin the maximum is attained at a
corner and evaluates in closed form:
\begin{equation}
\boxed{\; h(\theta) = a\,\lvert\cos\theta\rvert + b\,\lvert\sin\theta\rvert. \;}
\end{equation}
The absolute values appear because the best corner is whichever one lies on the
same side as $\mathbf{u}$.

The standoff radius on bearing $\theta$ is that reach plus a working clearance
$d_w$, floored at a minimum $r_{\min}$ so that a very small piece is not
approached too closely:
\begin{equation}
r(\theta) = \max\big(h(\theta) + d_w,\; r_{\min}\big),
\qquad d_w = r_{\min} = 0.8\ \text{m},
\end{equation}
and the candidate pose on that bearing is
\begin{equation}
\mathbf{x}(\theta) = \mathbf{c}_{xy} + r(\theta)\, \mathbf{u}(\theta),
\qquad
\psi(\theta) = \theta + \pi ,
\end{equation}
the yaw $\psi$ turning the robot back towards the centroid so the camera faces
the furniture. Bearings are sampled uniformly, $\theta_k = 2\pi k/n$ for
$k = 0,\dots,n-1$ with $n = 12$.
\textit{Implemented in} \texttt{core/navigation/standoff.py}, \texttt{ring\_radius} \textit{and} \texttt{candidates}.

The value $d_w = 0.8$~m is set by the navigation stack's inflation radius
($0.5$~m) and robot radius ($0.3$~m): closer than this and the goal lies in
inflated, high-cost space where the local controller struggles even when the
global planner accepts the goal.

\subsection{Rejecting poses inside other furniture}

The support function knows only about the target piece. A coffee table flanked by
sofas therefore produces top-ranked candidates inside a sofa. A candidate is
rejected if it lies within any other piece's footprint, padded by the robot's
radius $\rho = 0.3$~m:
\begin{equation}
\boxed{\;
\text{blocked}(\mathbf{x}) \iff
\exists j : \;
\lvert x - c^{j}_x \rvert \le \tfrac{D^{j}_x}{2} + \rho
\;\wedge\;
\lvert y - c^{j}_y \rvert \le \tfrac{D^{j}_y}{2} + \rho .
\;}
\end{equation}
\textit{Implemented in} \texttt{core/navigation/standoff.py}, \texttt{blocks}.

This is the standard axis-aligned box overlap test, with the robot's circular
footprint accounted for by growing the obstacle rather than the robot --- the
Minkowski-sum trick, exact here because a disc of radius $\rho$ swept around a
rectangle is contained in the rectangle grown by $\rho$ on each side.

\subsection{Ranking the candidates}

Surviving candidates are ordered lexicographically by the key
\begin{equation}
\kappa(\theta) = \Big( \big\lfloor r(\theta)/\delta \big\rceil,\;
\norm{\mathbf{x}(\theta) - \mathbf{x}_{\text{robot}}} \Big),
\qquad \delta = 0.1\ \text{m},
\end{equation}
where $\lfloor\cdot\rceil$ is rounding to the nearest integer. The first
component groups radii into $10$~cm bands and prefers standing close; the second
breaks ties by travel distance. Quantising first is what makes the tie-break
meaningful: without it, radii differing by a millimetre would never compare
equal and travel would never be consulted.

\subsection{Inverse reachability: where the arm can reach from}

Now the second question. Given a desired hand pose
$\mathbf{H} \in \SE(3)$ --- position $\mathbf{p}$, rotation $\mathbf{R}$ --- for
which base positions can the arm achieve it?

The arm has five actuated joints: a prismatic lift and four revolute joints. The
key structural fact, inherited from the analytic inverse kinematics, is that
three of the wrist axes intersect at a single point. Backing the palm offset off
the hand pose gives that intersection, the \emph{wrist centre}:
\begin{equation}
\boxed{\;
\mathbf{w} = \mathbf{p} + L_{81}\, \mathbf{R}\mathbf{e}_1 - L_{82}\, \mathbf{R}\mathbf{e}_3 .
\;}
\label{eq:wrist}
\end{equation}
\textit{Implemented in} \texttt{core/navigation/reach.py}, \texttt{\_wrist\_center}.

In words: start at the hand, step $L_{81}$ along the hand's own $x$ axis and
$L_{82}$ \emph{backwards} along its own $z$ axis. Because these are the hand's
axes and not the world's, the rotation $\mathbf{R}$ must be applied to them,
which is why the columns $\mathbf{R}\mathbf{e}_1$ and $\mathbf{R}\mathbf{e}_3$
appear. For the HSRC hand $L_{81} = 0$ and $L_{82} = 0.155$~m.

The wrist centre is decisive because the vertical prismatic lift can place it at
any height within its travel, independently of everything else. So:

\begin{itemize}
\item the \emph{height} $w_z$ alone decides whether the pose is reachable at all;
\item what remains in the horizontal plane is a \emph{circle} of base positions
      about $\mathbf{w}_{xy}$, since the base may stand anywhere at the right
      distance and turn to face it.
\end{itemize}

The set of admissible base positions is therefore an \emph{annulus} --- the
region between two concentric circles --- centred at $\mathbf{w}_{xy}$. This is
the closed-form answer to a question usually attacked by sampling a dense grid of
poses, solving inverse kinematics on each, and inverting the result into a
lookup table.

\subsection{Height feasibility}

Write $t_3 \in [t_3^{\min}, t_3^{\max}]$ for the lift extension, $L_3$ for the
lift's mounting height, and $L_{52}$ for the forearm length. The pose is
reachable at some base position only if
\begin{equation}
z_{\text{floor}} \;\le\; w_z \;\le\; t_3^{\max} + L_3 + L_{52},
\end{equation}
where the upper bound is the lift fully raised with the arm straight up, and
\begin{equation}
z_{\text{floor}} = t_3^{\min} + L_3 + L_{52}\cos t_4^{\min} + L_{51}\sin t_4^{\min}
\end{equation}
is the lift fully lowered with the arm flexed as far down as its limit
$t_4^{\min}$ allows. If $w_z$ falls outside this interval, \emph{no} amount of
driving helps, and the function returns ``unreachable'' rather than an annulus.
This is a cheap test worth running before committing to a drive.

\subsection{The annulus radii}

With the arm flexed by $t_4$, the horizontal distance from the base origin to the
wrist centre is
\begin{equation}
\boxed{\;
\rho(t_4) = \sqrt{\big(L_{52}\sin t_4 - L_{51}\cos t_4 - L_{41}\big)^2 + L_{42}^{\,2}} \;}
\label{eq:rho}
\end{equation}
\textit{Implemented in} \texttt{core/navigation/reach.py}, \texttt{\_radius}.

The first squared term is the reach in the arm's own plane of motion; the
constant $L_{42}$ is the lateral offset of the arm's mounting from the robot's
centreline, which contributes by Pythagoras because it is perpendicular to that
plane.

It is convenient to write the forearm geometry in polar form,
\begin{equation}
L = \sqrt{L_{51}^2 + L_{52}^2}, \qquad \alpha = \operatorname{atan2}(L_{52}, L_{51}),
\end{equation}
so that the wrist centre's height above the lift carriage is $L\sin(t_4 + \alpha)$.
Inverting this gives the flex angle that achieves a required height offset
$\Delta z$:
\begin{equation}
t_4(\Delta z) = \arcsin\!\left(\frac{\Delta z}{L}\right) - \alpha .
\end{equation}

\paragraph{Outer radius.} The arm reaches furthest horizontally when it is
straight out, $t_4 = -\alpha$. That is available only if the lift can supply the
required height on its own, i.e.\ $t_3^{\min} + L_3 \le w_z \le t_3^{\max} + L_3$.
Outside that band the arm must angle up or down to make up the difference, which
costs horizontal reach:
\begin{equation}
r_{\max} =
\begin{cases}
\rho\big(t_4(w_z - t_3^{\min} - L_3)\big), & w_z < t_3^{\min} + L_3,\\[4pt]
\rho(-\alpha), & t_3^{\min} + L_3 \le w_z \le t_3^{\max} + L_3,\\[4pt]
\rho\big(t_4(w_z - t_3^{\max} - L_3)\big), & w_z > t_3^{\max} + L_3 .
\end{cases}
\end{equation}

\paragraph{Inner radius.} Symmetrically, the smallest usable radius is set by how
far the arm can fold back and still hold the required height. With
$\phi = -t_4^{\min} - \alpha$ the fully folded configuration,
\begin{equation}
r_{\min} =
\begin{cases}
\sqrt{(L_{51} + L_{41})^2 + L_{42}^{\,2}}, & w_z \ge t_3^{\min} + L_3 + L_{52},\\[4pt]
\rho\big(t_4(w_z - t_3^{\min} - L_3)\big), & w_z > t_3^{\min} + L_3 + L_{52}\cos\phi + L_{51}\sin\phi,\\[4pt]
\rho(t_4^{\min}), & \text{otherwise.}
\end{cases}
\end{equation}

\subsection{Margin and the usable band}

Both edges of the annulus are \emph{singular} configurations: at $r_{\max}$ the
arm is fully extended, at $r_{\min}$ folded against the body. Near either, small
positioning errors produce large joint excursions, and the ability to make fine
corrections vanishes. Since the navigation goal tolerance is $0.25$~m --- enough
on its own to push a target off the edge --- a margin $\mu = 0.05$~m is kept
clear of both:
\begin{equation}
\mathcal{B} = [\, r_{\min} + \mu, \; r_{\max} - \mu \,],
\end{equation}
which is empty, and reported as such, when the ring is too thin. This is a
deliberate stand-in for the manipulability score a sampled reachability map would
have stored.

The test for ``can the arm reach this without driving'' is then simply
\begin{equation}
\boxed{\;
\text{reachable}(\mathbf{H}, \mathbf{x}) \iff
\norm{\mathbf{x} - \mathbf{w}_{xy}} \in \mathcal{B}.
\;}
\end{equation}
\textit{Implemented in} \texttt{core/navigation/reach.py}, \texttt{reachable\_from}.

\subsection{Base yaw}

The annulus gives \emph{where} to stand but not \emph{which way to face}. The
arm flexes only in its own plane and is mounted a distance $L_{42}$ off the
robot's centreline, so in the base frame the wrist centre always lies at exactly
\begin{equation}
{}^{B}\mathbf{w}_{xy} = \Big(\sqrt{r^2 - L_{42}^{\,2}},\; L_{42}\Big),
\qquad r = \norm{\mathbf{x} - \mathbf{w}_{xy}} .
\label{eq:invariant}
\end{equation}
That constraint pins the yaw. If $\beta$ is the world bearing from the base to
the wrist centre,
\begin{equation}
\beta = \operatorname{atan2}(w_y - y,\; w_x - x),
\end{equation}
then the required heading is
\begin{equation}
\boxed{\;
\psi = \beta - \arcsin\!\left(\frac{L_{42}}{r}\right).
\;}
\end{equation}
\textit{Implemented in} \texttt{core/navigation/reach.py}, \texttt{\_yaw}.

The $\arcsin$ term is the angle subtended by the lateral offset at distance $r$:
if the arm were mounted on the centreline ($L_{42}=0$) the robot would simply
face the wrist centre. Equation \eqref{eq:invariant} is exactly what the test
suite verifies to $10^{-9}$, and it is the sharpest check on the whole
derivation, since any yaw error shows up in it immediately.

\paragraph{A caution.} The arm has five degrees of freedom, while a general pose
in $\SE(3)$ has six. The missing degree of freedom is supplied by the base yaw
$\psi$. Consequently $\psi$ is not a preference but a constraint, and the
navigation stack's yaw tolerance of $0.25$~rad ($\approx 14^\circ$) is a direct
error on an arm with no redundancy to absorb it. Re-solving from the pose
actually achieved, rather than the one commanded, is the remedy.

\subsection{Choosing among the ring}

Every point of the usable band reaches the target equally well by construction,
so the tie-break is travel. Along each sampled bearing $\phi_k$ the code stands
at the radius within the band nearest to where the robot already is,
\begin{equation}
r_k = \operatorname{clip}\big(\norm{\mathbf{x}_{\text{robot}} - \mathbf{w}_{xy}},\;
r_{\min}+\mu,\; r_{\max}-\mu\big),
\end{equation}
and the resulting poses are sorted by
$\norm{\mathbf{x}_k - \mathbf{x}_{\text{robot}}}$. Which of them is
\emph{collision-free} is not decided here; that remains the planner's job.

\subsection{Seeing is not reaching}
\label{sec:disjoint}

The two rings never overlap. The standoff ring satisfies
$r(\theta) \ge d_w = 0.8$~m by construction, while the reach annulus satisfies
$r \le r_{\max} \le 0.492$~m for every attainable height. Hence
\begin{equation}
\boxed{\;
\min_\theta r(\theta) \;=\; 0.8\ \text{m} \;>\; 0.492\ \text{m} \;\ge\; r_{\max}.
\;}
\end{equation}
A pose from which the robot can see a piece of furniture is therefore
\emph{always} too far away to grasp anything on it. The second, shorter drive is
a mandatory stage of the pipeline rather than a refinement --- a fact asserted by
a test, so that if the constants ever change the consequence is a loud failure
rather than a silently redundant stage.

\subsection{Choosing between identical objects}

When several instances share a label, the one meant is the nearest to the robot,
measured in the plane:
\begin{equation}
i^\star = \arg\min_i \norm{\, \bar{\mathbf{p}}^{(i)}_{xy} - \mathbf{x}_{\text{robot}} \,}.
\end{equation}
Height is excluded because the robot drives on the floor: a can on a high shelf
is not further away for being high up. Absent a robot position, the most
confident instance is taken instead.
\textit{Implemented in} \texttt{core/scene\_graph/graph.py}, \texttt{find\_object}.

%======================================================================
\section{Constants}

All values are for the HSRC. The HSRB variant differs in the lift mount, palm
offset and lateral arm mount by millimetres that are nonetheless enough to miss
a grasp, so the robot is always named explicitly and never defaulted silently.

\begin{table}[h]
\centering
\begin{tabular}{llll}
\toprule
Symbol & Value & Units & Meaning \\
\midrule
$L_3$        & 0.350  & m & Lift mounting height \\
$L_{41}$     & 0.141  & m & Longitudinal arm mount offset \\
$L_{42}$     & 0.0785 & m & Lateral arm mount offset \\
$L_{51}$     & 0.005  & m & Forearm perpendicular offset \\
$L_{52}$     & 0.345  & m & Forearm length \\
$L_{81}$     & 0.0    & m & Palm offset along hand $x$ \\
$L_{82}$     & 0.155  & m & Palm offset along hand $z$ \\
$t_3$        & $[0,\,0.69]$    & m   & Lift travel \\
$t_4^{\min}$ & $-2.62$         & rad & Arm flex lower limit \\
$\mu$        & 0.05   & m & Margin kept clear of both annulus edges \\
$d_w$        & 0.8    & m & Standoff working distance \\
$\rho$       & 0.3    & m & Robot radius \\
$\delta$     & 0.1    & m & Standoff radius quantisation \\
\bottomrule
\end{tabular}
\end{table}

\begin{table}[h]
\centering
\begin{tabular}{ll}
\toprule
Derived quantity & Value \\
\midrule
Maximum horizontal reach & 0.492 m \\
Minimum reach (arm folded) & 0.166 m \\
Wrist-centre height envelope & 0.048 -- 1.385 m \\
Top-down grasp ceiling & 1.230 m \\
Navigation goal tolerance & 0.25 m, 0.25 rad \\
Octomap padding offset & 0.10 m \\
\bottomrule
\end{tabular}
\end{table}

Above the top-down ceiling of $1.230$~m a vertical approach is impossible and a
side approach is required; the palm offset works with the arm rather than against
it in that orientation, so the reachable height is greater.

\end{document}
